% Chapter 4

\chapter{实验与结果分析}

\section{实验设置}

\subsection{实验目标与验证思路}

本章围绕特定动画风格视频生成任务，重点回答四个问题：其一，数据整理与训练流程是否稳定可靠；其二，基于LoRA的参数高效微调能否在有限样本条件下学习到较为明确的风格特征；其三，不同底模先验是否会显著影响风格迁移的方向与强度；其四，图像阶段获得的风格能力能否在短视频原型中延续，并在一定程度上改善跨帧一致性。围绕上述目标，实验设计采用“短训验证—长训加强—底模替换—视频原型扩展”的渐进路线，既保证每一阶段都有明确结论，又能使后续实验建立在前一阶段可复现的结果之上。

图像阶段首先以Stable Diffusion 1.5为基础模型完成300步短训实验，用于验证768分辨率训练、LoRA注入和批量推理链路是否能够稳定运行。随后，在保持主体流程不变的条件下，对caption的组织方式和训练步数进行扩展，观察风格强化策略是否能够带来更加明显的手绘动画表现。考虑到风格生成任务不仅受训练数据影响，还明显受到底模视觉先验的制约，本文进一步引入更偏故事书插画风格的StorybookRedmond以及具有较强二次元风格先验的Counterfeit-V3.0作为替代底模，以此考察不同底模对目标风格适配效果、结构保真度和生成可控性的影响。

在实验方法上，所有图像对比均尽量控制提示词、随机种子、采样步数和输出分辨率一致，并通过统一脚本完成baseline与LoRA结果的并行生成。这样做的目的，是尽量将观察到的差异归因于LoRA权重与底模组合本身，而不是推理噪声或人工挑选样例带来的偶然偏差。视频阶段则延续先图像、后视频的思路，不直接训练大规模视频扩散模型，而是在图像风格LoRA的基础上构建递推式短视频生成原型，通过无参考时序一致性指标评估相邻帧变化是否更平滑、边缘结构是否更稳定，以及亮度闪烁是否得到抑制。

\subsection{数据集与预处理结果}

图像风格学习阶段采用整理后的动画风格图像数据集作为训练基础。经过扫描与预处理后，共得到810张样本，其中有效样本810张，无损坏或无法读取样本；按照9:1的比例划分后，训练集为729张，验证集为81张。所有样本均配有同名文本描述文件，因此caption覆盖率达到100\%，这为后续的文本条件训练提供了较为一致的输入基础。由于本研究并不追求大规模通用模型训练，而是强调特定风格的定向学习，因此数据集规模被控制在中等水平，使单轮实验能够在有限时间内完整闭环。

在预处理过程中，全部图像统一转换为RGB格式，并执行中心裁剪、比例规整和尺寸缩放，以消除原始采集源中普遍存在的宽高差异和边缘留白问题。实验分辨率统一设置为768×768，这一设置兼顾了画面细节、显存占用与训练效率。

视频阶段的数据主要服务于流程验证与一致性分析，而非大规模视频底模训练。当前共整理出5个片段样本，其中训练片段4个、验证片段1个，累计导出帧数100帧。进一步的交叉检查表明，训练集与验证集片段之间不存在编号重叠，数据泄漏计数为0，可认为划分过程满足基本的独立性要求。另一方面，图像推理阶段还预先定义了10条场景提示词，覆盖港口、街道、铁路、桥梁、卧室、山村、森林小径和河边集市等典型动画场景，为统一的横向比较提供了稳定测试集合。

\begin{table}[htbp]
  \centering
  \caption{数据集预处理与划分结果}
  \makebox[\textwidth][c]{%
  \begin{tabularx}{\whutablewidth}{CCCCCCCC}
    \toprule
    数据类型 & 扫描数量 & 有效数量 & 训练集 & 验证集 & \makecell[c]{caption\\数量} & \makecell[c]{缺失\\caption} & 说明 \\
    \midrule
    \makecell[c]{图像风格\\数据集} & 810 & 810 & 729 & 81 & 810 & 0 & \makecell[c]{统一处理为\\768$\times$768} \\
    \makecell[c]{视频片段\\数据集} & 5 & 5 & 4 & 1 & 5 & 0 & \makecell[c]{累计导出\\100 帧，无\\数据泄漏} \\
    \bottomrule
  \end{tabularx}%
  }
\end{table}

\subsection{实验环境与参数配置}

实验的训练环境采用基于Conda管理的独立虚拟环境，以避免不同模型依赖冲突对实验结果造成干扰，整个实验过程均在Linux系统下运行，训练框架以PyTorch、Diffusers、Transformers、Accelerate和PEFT为核心，形成相对稳定的扩散模型微调工作流。

\begin{table}[htbp]
  \centering
  \caption{实验硬件与软件环境}
  \makebox[\textwidth][c]{%
  \begin{tabularx}{\whutablewidth}{
    >{\hsize=0.75\hsize\centering\arraybackslash}X
    >{\hsize=1.00\hsize\centering\arraybackslash}X
    >{\hsize=1.25\hsize\centering\arraybackslash}X
  }
    \toprule
    项目 & 配置 & 说明 \\
    \midrule
    GPU & \makecell[c]{NVIDIA RTX 3080\\Laptop} & \makecell[c]{承担 768 分辨率 LoRA 训练与批量\\推理} \\
    操作系统 & Linux 5.15.0-78-generic & 实验脚本运行环境 \\
    Python & 3.10.8 & 核心运行时 \\
    PyTorch & 2.11.0+cu130 & 深度学习训练框架 \\
    Diffusers & 0.27.2 & 扩散模型加载与推理框架 \\
    Transformers & 4.39.3 & 文本编码与模型组件支持 \\
    Accelerate & 0.28.0 & 训练加速与混合精度支持 \\
    PEFT & 0.10.0 & LoRA参数高效微调支持 \\
    \bottomrule
  \end{tabularx}%
  }
\end{table}

训练参数的主体设置在多轮实验中保持大致一致：图像分辨率固定为768，训练批量大小为1，梯度累积步数为2，学习率统一为1e-04，LoRA秩设为16，并启用fp16 混合精度与TF32加速。由于目标在于对比不同底模和训练策略对风格学习的影响，而非追求某一模型的极限拟合，因此参数设计强调公平与可重复：Stable Diffusion 1.5主线实验允许更长训练步数，StorybookRedmond和Counterfeit-V3.0则采用相同的 1200步设置，以便观察不同底模先验在相近训练强度下的表现差异。

\begin{table}[htbp]
  \centering
  \caption{主要训练方案配置}
  \makebox[\textwidth][c]{%
  \begin{tabularx}{\whutablewidth}{
    >{\hsize=1.8\hsize\centering\arraybackslash}X
    >{\hsize=2.1\hsize\centering\arraybackslash}X
    >{\hsize=0.8\hsize\centering\arraybackslash}X
    >{\hsize=0.6\hsize\centering\arraybackslash}X
    >{\hsize=0.5\hsize\centering\arraybackslash}X
    >{\hsize=0.6\hsize\centering\arraybackslash}X
    >{\hsize=0.6\hsize\centering\arraybackslash}X
  }
    \toprule
    实验名称 & 底模 & \makecell[c]{最大训\\练步数} & 分辨率 & \makecell[c]{LoRA\\秩} & 学习率 & \makecell[c]{caption\\dropout} \\
    \midrule
    SD1.5 短训 & \makecell[c]{runwayml/stable-diffusion-\\v1-5} & 300 & 768 & 16 & 1e-4 & 0.05 \\
    \makecell[c]{caption 强\\化短训} & \makecell[c]{runwayml/stable-diffusion-\\v1-5} & 300 & 768 & 16 & 1e-4 & 0.05 \\
    \makecell[c]{caption 强\\化长训} & \makecell[c]{runwayml/stable-diffusion-\\v1-5} & 2000 & 768 & 16 & 1e-4 & 0.05 \\
    \makecell[c]{StorybookRedmond 长\\训} & Yntec/StorybookRedmond & 1200 & 768 & 16 & 1e-4 & 0.00 \\
    \makecell[c]{Counterfeit-V3.0 长训} & danbrown/Counterfeit-v3-0 & 1200 & 768 & 16 & 1e-4 & 0.00 \\
    \bottomrule
  \end{tabularx}%
  }
\end{table}

推理阶段同样遵循控制变量原则。单图与批量对比均固定随机种子和场景提示词，仅在是否加载LoRA、使用何种底模以及采样参数细节上进行调整。总体上，Stable Diffusion 1.5主线实验的推理步数略高于替代底模，以保证其在较中性的底模先验下仍能获得充分去噪时间；StorybookRedmond和Counterfeit-V3.0则适当降低采样步数与guidance scale，以减少过强引导造成的图像硬化与局部失真。视频原型实验在此基础上进一步引入固定帧数、强度参数和帧率设置，使不同组视频的一致性指标具备可横向比较的意义。

\subsection{评价指标与结果呈现方式}

为了避免仅凭主观印象判断风格学习效果，本文将训练过程指标、图像对比指标和视频一致性指标结合起来进行综合分析。训练阶段记录global step与best loss，用于反映LoRA适配器在去噪目标上的收敛情况；但需要指出，扩散模型训练损失只反映噪声预测误差，并不能直接等价于“更接近吉卜力风格”或“视觉效果更好”。因此，本章在解释loss变化时始终将其与后续的图像结果、差异图和定量对比表结合，避免以单一训练曲线替代最终质量判断。

图像阶段的核心定量指标包括像素平均绝对差、亮度变化、对比度变化、饱和度变化和边缘强度变化。像素平均绝对差用于衡量加载LoRA后图像整体发生了多大程度的变化，其数值越大，通常说明LoRA对底模输出的干预越强；但这一指标只能说明变化幅度，并不能自动说明这种变化一定朝着目标风格优化。亮度、对比度和饱和度等指标则有助于进一步分析风格偏移的方向，例如更柔和的手绘动画风格往往伴随色彩压缩和局部边缘处理的变化，而边缘强度变化可以辅助判断画面轮廓是否趋向平滑。

视频阶段使用无参考的一致性指标来衡量跨帧稳定性。MAD描述相邻帧像素差异的平均水平，数值越低通常表示帧间过渡更平滑；MED描述相邻帧边缘结构差异，数值越低通常表示结构变化更稳定；LS描述全局亮度均值的标准差，数值越低通常表示亮度闪烁更弱。由于本研究的视频部分仍处于原型验证阶段，尚未引入FVD等更复杂的时空评测指标，因此MAD、MED与LS的作用更多是提供趋势性证据，用以判断图像阶段训练得到的风格LoRA是否对短视频连续性具有正向帮助。

\section{图像风格学习实验结果分析}

\subsection{基于Stable Diffusion 1.5的短训验证结果}

首先考察基于Stable Diffusion 1.5的300步短训结果。该实验的best loss为0.001998，像素平均绝对差为0.111979。从数值上看，LoRA已经能够对底模输出产生可观察的风格干预，但干预程度仍保持在相对温和的范围内。这种结果符合短训实验的预期定位：其核心目标不是追求最强烈的风格迁移，而是确认在新的云端算力平台上，数据预处理、LoRA注入、训练、推理和对比可视化是否构成稳定闭环。

从汇总指标看，该轮实验的亮度变化为0.0152，对比度变化为-0.0090，饱和度变化为-0.0537，边缘强度变化为-0.0070。这些指标共同表明，LoRA并未简单地把图像推向极端高饱和或极端高对比方向，而是在颜色分布、线条强度和局部结构上进行较为克制的重塑。从工程角度看，这种变化有限但方向明确的状态说明训练链路是有效的，同时也提示仅靠300步训练尚不足以充分放大目标风格特征。

\begin{figure}[htbp]
  \centering
  \includegraphics[width=0.92\textwidth]{chapter4_figure1_1.png}
  \caption{基于Stable Diffusion 1.5短训结果的baseline与LoRA对比样例}
\end{figure}

从典型样例可以看到，加载LoRA后的图像在构图、建筑分布和前景元素上与 baseline存在明显差异，说明模型已经不再只是微调局部色彩，而是开始对场景组织方式产生影响。与此同时，图像整体仍保持了较高的可读性和场景完整性，没有出现大范围语义崩坏。这意味着短训LoRA在本任务中具备较好的冒烟验证价值：它能够迅速证明风格学习是可行的，并为后续更长训练和更换底模提供可靠起点。

进一步观察如图4-2所示的10条提示词的差异分布，可以发现不同场景对风格LoRA的敏感度并不一致。例如，包含复杂建筑群、前景栏杆或交通元素的提示词往往产生更大的像素差异，这说明风格学习并不是均匀作用于整张图像，而会优先改变更能承载手绘笔触与空间层次的区域。对于论文写作而言，这一结果具有两个意义：一方面，它证明LoRA已经具备一定的场景适应能力；另一方面，它也提示后续实验不能只看单一样例，而应结合多场景统计结果来判断风格学习是否稳定。

\begin{figure}[htbp]
  \centering
  \includegraphics[width=0.92\textwidth]{chapter4_figure2_1.png}
  \caption{Stable Diffusion 1.5短训实验在10条提示词上的差异指标分布}
\end{figure}

\subsection{caption组织增强实验分析}

在完成基础短训验证后，本文尝试通过轻量化的caption风格增强策略进一步放大目标风格信号。其出发点是，若文本描述能够在内容信息之外加入更统一的风格表述，模型可能会更快学习到手绘动画的色调、背景表现和镜头氛围。然而实验结果表明，这一策略并未稳定带来正向收益。300步caption强化短训的像素平均绝对差为0.106736，略低于基础短训的0.111979，说明在短步数条件下，额外的caption风格表述并没有显著增强风格注入强度，反而可能在一定程度上削弱了底模对场景语义的自由组织。

当训练步数增加到2000步后，caption强化长训的best loss下降到0.001930，像素平均绝对差上升到0.147990，相较于基础短训提高了32.16\%。如果只看训练损失和整体变化幅度，这一结果似乎意味着模型学到了更强的风格偏移能力；但结合实际样例与单提示词统计可以发现，变化增强并不等价于质量提升。部分场景在长训后出现了更明显的构图重排、主体位置漂移和局部细节失衡，尤其是包含轨道、车体、远景透视线或复杂建筑排列的提示词，更容易出现与原始语义布局不完全一致的情况。

这一现象说明，过度依赖统一化caption风格描述可能使模型在训练中更倾向于追逐风格标签的表面特征，而弱化对具体场景语义的细致响应。对于扩散模型而言，caption本质上是条件约束的一部分，若文本中过多重复统一风格短语，而样本规模又有限，模型容易将这些高频词与某些固定色块、构图模板或光照模式绑定，最终在推理阶段表现为风格更强但可控性下降。换言之，caption强化实验提供了一条重要的负面经验：在样本规模不大的前提下，训练目标应优先保证内容-风格之间的平衡，而不是单纯追求风格词堆叠。

因此，caption组织增强实验虽然在数值上带来了更大的风格差异，但从所关注的接近目标风格且保持场景可读性的综合目标出发，它并未成为最优方案。该实验的意义主要在于验证了一点：针对特定风格的LoRA训练，文本描述质量固然重要，但统一强化并非越强越好，过度一致的文本模板有可能压缩样本本身的内容多样性，进而导致生成结果趋向模板化。

\subsection{不同底模长训结果比较}

在确认主线流程可用之后，本文进一步将同一数据集接入不同底模，考察底模先验对目标风格学习的影响。如表4-4所示，比较对象包括Stable Diffusion 1.5主线结果、StorybookRedmond长训结果和Counterfeit-V3.0长训结果。三者的best loss分别为0.001998、0.001994和0.001902。如果从loss角度看，Counterfeit-V3.0的收敛值最低，似乎更容易被数据集驱动；但结合图像对比指标可以看到，底模之间的差异并不能只简单考虑是否更容易收敛。

\begin{table}[htbp]
  \centering
  \caption{各组图像训练结果统计}
  \makebox[\textwidth][c]{%
  \begin{tabularx}{\whutablewidth}{
    >{\hsize=1.80\hsize\centering\arraybackslash}X
    >{\hsize=0.75\hsize\centering\arraybackslash}X
    >{\hsize=1.00\hsize\centering\arraybackslash}X
    >{\hsize=1.10\hsize\centering\arraybackslash}X
    >{\hsize=0.75\hsize\centering\arraybackslash}X
    >{\hsize=0.85\hsize\centering\arraybackslash}X
    >{\hsize=0.85\hsize\centering\arraybackslash}X
    >{\hsize=0.90\hsize\centering\arraybackslash}X
  }
    \toprule
    实验名称 & \makecell[c]{global\\step} & best loss & \makecell[c]{像素平均\\绝对差} & \makecell[c]{亮度变\\化} & \makecell[c]{对比度\\变化} & \makecell[c]{饱和度\\变化} & \makecell[c]{边缘强\\度变化} \\
    \midrule
    SD1.5 短训 & 300 & 0.001998 & 0.111979 & 0.0152 & -0.0090 & -0.0537 & -0.0070 \\
    caption 强化长训 & 2000 & 0.001930 & 0.147990 & 0.0134 & -0.0012 & -0.0150 & -0.0088 \\
    \makecell[c]{StorybookRedmond\\长训} & 1200 & 0.001994 & 0.197670 & 0.0047 & -0.0117 & -0.0898 & -0.0001 \\
    \makecell[c]{Counterfeit-V3.0\\长训} & 1200 & 0.001902 & 0.225164 & 0.0090 & -0.0288 & -0.0241 & 0.0055 \\
    \bottomrule
  \end{tabularx}%
  }
\end{table}

从整体变化幅度看，StorybookRedmon长训的像素平均绝对差达到0.197670，相较于Stable Diffusion 1.5短训提高了76.52\%；Counterfeit-V3.0长训则达到0.225164，提升幅度达到101.08\%。这说明具有更强动画插画或二次元先验的底模，的确更容易在少量LoRA微调之后形成明显的风格偏移。对于追求风格变化强度的实验目标而言，这样的底模更有优势，因为它们本身就处于更接近目标分布的视觉空间中，LoRA只需在原有先验上进行定向修正即可。

\begin{figure}[htbp]
  \centering
  \includegraphics[width=0.92\textwidth]{chapter4_figure3_1.png}
  \caption{StorybookRedmond长训结果的baseline与LoRA对比样例}
\end{figure}

但从具体样例可以进一步看到，底模先验的增强也会带来新的问题。StorybookRedmond在若干场景中能够更快形成柔和、插画感较强的色面组织，使画面更接近童话绘本或动画海报的视觉风格；然而它也更容易重写原始构图关系，例如在海港、街巷等需要明显透视关系的场景中，局部建筑朝向、远景层次和水面空间感有时会被简化为更平面化的组织方式。也就是说，StorybookRedmond的风格表达更积极，但在维持提示词内容的几何可解释性方面并不一定优于更中性的Stable Diffusion 1.5。

\begin{figure}[htbp]
  \centering
  \includegraphics[width=0.92\textwidth]{chapter4_figure4_1.png}
  \caption{Counterfeit-V3.0长训结果的baseline与LoRA对比样例}
\end{figure}

Counterfeit-V3.0的表现则更加鲜明。该底模在轮廓线、色块边界和二次元角色化视觉语言上具有更强先验，因此加载LoRA后的变化最为剧烈，部分样例呈现出明显的动漫化重绘倾向。其优点在于风格冲击力强、轮廓表达清晰，适合快速生成具有动画感的画面；缺点则在于目标风格会被更广义的二次元审美所稀释，导致结果更接近泛动漫插画，而不一定稳定落在“温润、背景绘画感较强”的吉卜力式区域。尤其在港口、河道和街区等场景中，Counterfeit-V3.0更容易出现大尺度构图重排和结构关系改写，这在一定程度上削弱了提示词与图像之间的对应精度。

综合来看，不同底模之间不存在绝对单一的最优解，而是体现出明显的目标差异。若实验关注的是风格偏移强度，StorybookRedmond与Counterfeit-V3.0的长训结果显然比Stable Diffusion 1.5更激进；若实验关注的是风格注入与场景保真之间的平衡，则Stable Diffusion 1.5虽然变化幅度较小，却保留了更高的可控性和更稳定的场景语义响应。

\subsection{图像阶段结果归纳}

将以上实验放在同一框架下可以得到较为清晰的结论。首先，基于810张风格图像、768分辨率和LoRA参数高效微调的训练路线是可行的，图像实验已经从能否跑通进入如何平衡风格强度与内容保真的阶段。其次，训练损失的下降和像素差异的增大虽然能够说明模型对数据集的适配程度在上升，但这并不自动代表结果更接近目标风格，仍需要结合场景可读性和主观观感综合判断。再次，caption统一强化和更换强先验底模都可以放大风格迁移强度，但前者容易造成模板化，后者容易引入底模固有审美偏差，因此都需要在可控性和目标风格纯度之间做权衡。

\section{视频原型实验与时序一致性分析}

\subsection{视频原型实验方案}

考虑到专门视频扩散模型训练的成本远高于图像LoRA微调，本文在视频阶段采用递推式短视频生成原型进行验证。具体做法是先依据文本提示词生成首帧，再将上一帧作为下一帧的输入，通过固定strength、采样步数和随机种子递推得到一段短视频片段。该方案不需要额外训练庞大的视频主干网络，能够在已有图像风格LoRA的基础上，直接观察手绘动画风格是否能够延续到连续帧中，因此非常适合当前阶段的原型验证路线。

实验中首先选取10条具有代表性的场景提示词用于统一的图像对比，其中5条提示词被进一步用于12帧短视频的批量生成评估；在此基础上，又额外抽取2条提示词进行云端复核，以验证结果趋势是否具有稳定性。视频生成的主要参数包括固定帧数、固定种子基数、固定guidance scale和固定强度系数，保证baseline视频与加载LoRA后的视频具有一一对应关系。这样，指标的变化就可以更明确地反映LoRA对跨帧连续性的影响，而不是采样配置变化造成的噪声。

从数据组织角度看，视频预处理阶段共整理出5个可用片段并导出100帧样本，这为后续设计训练-验证片段划分、检查数据泄漏以及实现基本一致性度量提供了基础。尽管现阶段的视频样本规模较小，不足以支撑完整的时空生成模型训练，但它已经足以作为图像风格能力向连续帧迁移的证据来源。

\subsection{时序一致性定量结果}

\begin{table}[htbp]
  \centering
  \caption{短视频原型时序一致性指标对比}
  \makebox[\textwidth][c]{%
  \begin{tabularx}{\whutablewidth}{CCCCCCCCCC}
    \toprule
    \makecell[c]{评估\\集合} & \makecell[c]{Baseline\\MAD} & \makecell[c]{LoRA\\MAD} & \makecell[c]{MAD\\变化} & \makecell[c]{Baseline\\MED} & \makecell[c]{LoRA\\MED} & \makecell[c]{MED\\变化} & \makecell[c]{Baseline\\LS} & LoRA LS & LS 变化 \\
    \midrule
    \makecell[c]{5 组\\批量\\实验} & 0.13879 & 0.12986 & -6.43\% & 0.06252 & 0.06086 & -2.65\% & 0.01675 & 0.01872 & +11.72\% \\
    \makecell[c]{2 组\\复核\\实验} & 0.13898 & 0.13642 & -1.84\% & 0.06290 & 0.05993 & -4.71\% & 0.00817 & 0.00880 & +7.77\% \\
    \bottomrule
  \end{tabularx}%
  }
\end{table}

从5组批量实验结果看，加载LoRA后的视频在MAD指标上由0.138789下降到0.129858，降幅约为6.43\%；MED指标由0.062519下降到0.060861，降幅约为2.65\%。这说明在当前递推式视频生成方案下，图像阶段训练得到的风格LoRA确实能够在一定程度上降低相邻帧之间的整体变化与边缘结构突变，使画面过渡更平滑，结构抖动有所缓解。

但在亮度稳定性方面，LS指标从0.016752上升到0.018716，增幅约为11.72\%。也就是说，LoRA版本虽然在画面连续性和结构稳定性上略有改善，却没有同步抑制全局亮度闪烁，反而使亮度均值波动有所增加。结合生成机制可以推测，这是因为递推式img2img方案在每一帧都会重新执行带噪去噪过程，LoRA对色彩与光照风格的强化会在帧间逐步积累，从而把原本较弱的亮度变化放大为可测量的全局波动。

2组复核实验进一步验证了这一趋势。复核结果中，LoRA版本的MAD由0.138983降至0.136421，MED由0.062896降至0.059934，对应降幅分别为1.84\%和4.71\%；与此同时，LS仍从0.008172上升到0.008806。尽管复核样本数量较少，但与批量实验呈现出相同方向的变化，这说明结构连续性改善而亮度稳定性不足的结论具有一定稳定性，而不是个别样例造成的偶然现象。

从指标解释的角度看，这组结果也符合当前原型方案的技术特征。递推式生成的优势在于，它天然以上一帧为条件，因此能够较好保留颜色分区、主体轮廓和背景布局的局部延续性；但其不足在于，它缺乏真正的视频时序建模模块，无法显式学习全局运动轨迹、镜头规律和长期亮度约束。因此，只要连续帧中的局部扰动不断累积，就可能在颜色层面产生轻微漂移或整体曝光起伏。图像风格能力已经开始向连续帧迁移，但距离成熟的视频扩散系统仍有明显差距。

\subsection{视频结果的定性观察与局限}

结合实际生成的短视频片段可以进一步观察到，LoRA对连续帧的作用主要体现在前几帧较为明显，表现为画面整体色调、天空与建筑色块的组织方式更趋统一，局部笔触和边缘线条也更符合目标风格。但随着递推步数增加，误差会逐帧积累，某些片段会出现局部区域变形、纹理重复、结构轻微漂移甚至主体边界模糊等现象。尤其是在场景中存在水面反光、栏杆、轨道、电线杆或密集屋顶时，这类高频结构对递推生成更为敏感，容易被不断重绘而产生形态抖动。

这说明图像阶段训练得到的LoRA权重虽然能够提供风格方向上的约束，但并不能替代显式的时序建模。换句话说，LoRA解决的是每一帧更像目标风格的问题，而视频生成还需要额外解决每一帧之间如何保持运动与结构连续的问题。当前实验已经证明，单帧风格能力向视频序列迁移是可行的，但要想进一步提升亮度稳定性、主体一致性和长序列可控性，仍需要引入专门面向视频的运动先验、光流约束或跨帧特征交互机制。

尽管如此，视频原型实验仍然具有明确研究价值。首先，它使本文中的技术路线形成了完整闭环，即从图像数据构建、LoRA风格学习、图像推理对比，一直延伸到连续帧的原型验证；其次，它给出了定量证据，说明图像LoRA对短视频的结构平滑性存在一定正向影响；最后，它准确暴露出当前方案的瓶颈所在，为后续改进提供了清晰方向。

\section{实验讨论}

\subsection{风格学习有效性讨论}

综合图像与视频实验，可以确认本研究所采用的“中等规模风格图像数据集 + Stable Diffusion系列底模 + LoRA参数高效微调 + 批量对比评估”的总体方案是有效的。数据预处理阶段保证了图像与文本描述的一一对应关系，降低了训练条件噪声；LoRA训练阶段以较小参数量完成了风格迁移；批量推理与可视化模块则为结果分析提供了系统化支撑，使实验具备了多场景、多指标的综合比较依据。

实验还表明，目标风格学习的关键并不只在于训练步数是否足够长，更在于数据、文本和底模先验三者之间是否匹配。当底模与目标风格距离较远时，即使训练过程稳定，LoRA也可能只能形成有限的风格修正\cite{ref40}；当底模本身带有较强的动画先验时，模型又可能在风格方向上走得过远，把结果拉向更宽泛的二次元视觉分布。这一结论说明，在特定风格生成任务中，底模选择本身就是实验设计的重要变量，而不应被视为可忽略的背景条件。

\subsection{指标与主观观感差异的原因分析}

本章实验中一个较为典型的现象，是部分实验在训练损失下降或像素差异增大的同时，主观观感并未同步提升。造成这一现象的原因主要有三点。第一，扩散模型训练的优化目标是噪声预测误差，而不是直接优化风格相似度；因此loss的下降只能说明模型更适配当前训练数据的统计分布，无法直接断言生成结果一定更美观或更贴近目标风格。第二，像素平均绝对差等指标衡量的是变化幅度，而非变化质量，过大的变化既可能代表成功的风格注入，也可能意味着语义偏移、构图重写或局部失真。第三，风格任务本身具有较强审美属性，仍然需要结合图像样例进行人工判断。

以caption强化长训为例，其loss与差异幅度都优于基础短训，但实际效果并没有形成更稳健的目标风格，反而在部分提示词上出现了过度强化和模板化倾向。再以Counterfeit-V3.0为例，虽然它在数值上呈现出最强烈的风格偏移，但这种偏移并不完全等价于吉卜力式风格增强，因为模型同时引入了更强的泛动漫先验。由此可以得到一个更具普遍性的认识：在风格生成实验中，定量指标适合用于筛查趋势、定位问题和支撑比较，但最终判断仍应建立在定量分析、可视化样例任务、目标一致性三者结合的基础上。

\subsection{当前方案的局限与后续优化方向}

当前实验仍然存在几方面局限。其一，图像数据集规模虽然已经能够支撑LoRA风格学习，但样本覆盖面仍然有限，特别是在人物、复杂镜头和强运动场景上分布不足，这会限制模型对更丰富场景的泛化能力。其二，视频阶段尚未采用真正意义上的直接训练大规模视频扩散模型，因此跨帧一致性的提升幅度有限，亮度闪烁问题也未得到根本解决。其三，评价体系仍以无参考指标和人工观察为主，缺少更系统的图文一致性、感知质量和时序质量评价指标。

基于以上局限，后续优化可以从三个方向展开。首先，在数据层面继续扩充风格样本，并对caption做更细粒度的语义约束，而不是简单叠加统一风格词汇，以提高内容保持与风格强化的平衡性。其次，在模型层面引入更适合视频任务的结构，如AnimateDiff类运动模块、跨帧注意力机制或ControlNet条件约束，以降低长序列递推中出现的结构漂移和亮度抖动。再次，在评测层面增加CLIP相似度、感知距离以及更成熟的视频一致性指标，使结果分析更加全面。

总体而言，当前研究已经完成了从图像风格学习到短视频原型验证的关键链路搭建，并在多轮实验中积累了可直接写入论文的数据与现象。虽然距离完整、高质量的视频生成系统仍有差距，但本章实验已经足以说明所提出的技术路线具备可行性，并为后续继续完善模型、数据和评测体系奠定了明确基础。

\section{本章小结}

本章围绕实验设置、图像风格学习结果、短视频原型实验和综合讨论，对系统进行了较为完整的实验性验证。结果表明：在810张风格图像和768分辨率设置下，基于LoRA的参数高效微调能够有效学习目标风格特征，并在多场景提示词上产生稳定可观测的风格偏移；不同底模对结果影响显著，其中StorybookRedmond与Counterfeit-V3.0可带来更强的风格迁移，但也更容易引入构图重写和底模先验偏差；递推式视频原型实验显示，加载LoRA后的短视频在帧间平滑性和结构稳定性上有一定改善，但亮度闪烁仍需进一步优化。上述结论既验证了本文技术路线的可行性，也为下一章的总结与展望提供了实验依据。
